\begin{textsample}{POS Dim 1 – 2022\_c6 – Score -2.00 – 2022\_c6/enade\_2022\_tec\_proc\_ger\_587bf59e0a.txt}  \label{ex:f1_pos_264}
CAPÍTULO 7 ANÁLISE TÉCNICA DA PROVA Este capítulo tem por objetivo apresentar o desempenho dos estudantes concluintes de Tecnologia em Processos Gerenciais no Enade/2022.
Para isso, foram calculadas as estatísticas básicas da prova em seu todo ( Subseção 7.1.1 ), bem como as estatísticas relacionadas aos Componentes de Formação Geral ( Subseção 7.1.2 ) e de Conhecimento Específico da Área ( Subseção 7.1.3 ).
Dadas as suas características, foram analisadas, em separado, as questões objetivas ( Seção 7.2 ) e as questões discursivas ( Seção 7.3 ).
Tomando -se como base as duas questões discursivas do Componente de Formação Geral, nas Subseções 7.3.1.6 a 7.3.1.8, são apresentados os resultados e os comentários sobre a correção das respostas em relação à Língua Portuguesa.
Para as questões objetivas, foram disponibilizados os índices de Facilidade e de Discriminação Ponto-Bisserial, também em separado, para os Componentes de Formação Geral ( Subseção 7.2.1 ) e de Conhecimento Específico ( Subseção 7.2.2 ).
De cada componente, uma das questões foi escolhida para exemplificar a análise gráfica, relacionando -se as alternativas escolhidas pelos estudantes ( inclusive o gabarito ) com o número de acertos no componente.
No Anexo IV, apresenta -se a íntegra da análise gráfica para todas as questões objetivas.
Para cada uma das questões discursivas, os conteúdos dos tipos mais comuns de respostas dos estudantes são apresentados e comparados com o padrão de resposta esperado ( ver Anexo IX ).
Nas tabelas, constam as seguintes estatísticas das notas10 : média do desempenho na prova, erro-padrão da média, desvio padrão, nota mínima, mediana e nota máxima para cada um de seus componentes.
Tais estatísticas contemplam o total de estudantes concluintes da Área de Tecnologia em Processos Gerenciais inscritos e presentes na prova do Enade/2022, tendo em vista agregações, ou por grandes regiões e o país como um todo, ou por categoria administrativa, organização acadêmica da IES e modalidade de oferta dos cursos.
Em relação aos gráficos de distribuição de notas, o intervalo considerado foi de 10 unidades, aberto à esquerda e fechado à direita, com exceção do primeiro intervalo, [ 0 ; 10 ], fechado em ambos os 10 Uma definição dessas estatísticas pode ser encontrada no Glossário inserido no final desse relatório. 145 7.1 ESTATÍSTICAS BÁSICAS DA PROVA Esta seção apresenta estatísticas selecionadas e histogramas da nota geral ( 7.1.1 ) e de cada componente : Formação Geral ( 7.1.2 ) e Conhecimento Específico ( 7.1.3 ).
São também apresentadas estatísticas selecionadas de subpopulações, caracterizadas por grande região, categoria administrativa, organização acadêmica e modalidade de oferta. 7.1.1 ESTATÍSTICAS BÁSICAS GERAIS Na Tabela 7.1, são apresentadas as Estatísticas básicas da prova, por grande região, dos estudantes concluintes de Tecnologia em Processos Gerenciais.
Como se pode verificar a partir dos dados mostrados na Tabela 7.1, a média das notas da prova, como um todo, foi de 43,9 ( nas seções seguintes, serão analisados os Componentes de Formação Geral e de Conhecimento Específico ).
Os estudantes da região Norte obtiveram a média mais baixa ( 41,1 ), e os da região Sudeste alcançaram a média mais alta ( 44,7 ).
As demais Médias foram : 43,6 na região Sul, 42,8 na região Centro-Oeste e 41,7 na região Nordeste.
O desvio-padrão para o Brasil foi de 14,1, sendo o maior desvio-padrão encontrado na região Sudeste ( 14,4 ), e o menor, na região Norte ( 13,3 ), indicando uma dispersão um pouco menor das notas desta última região.
A região que obteve a maior nota Máxima foi a região Centro-Oeste ( 92,2 ), ao passo que a região que atingiu a menor nota Máxima foi a Norte ( 73,9 ).
A Mediana do Brasil, como um todo, foi de 43,5, sendo a maior Mediana obtida na região Sudeste ( 44,6 ), e a menor obtida na região Norte ( 40,9 ).
As regiões Sul e Sudeste foram as únicas que obtiveram nota mínima zero.
O comportamento das notas dos estudantes de todo o Brasil pode ser observado no Gráfico 7.1, que apresenta um histograma com a distribuição destas : uma distribuição unimodal com a moda no intervalo ( 40 ; 50 ].
Na Tabela 7.2, são apresentadas informações referentes à Média da nota final, desagregadas por categoria administrativa, organização acadêmica e modalidade de oferta dos cursos.
Os dados da Tabela 7.2 mostram que, em relação à categoria administrativa, os estudantes das IES públicas obtiveram Média mais alta ( 51,1 ) que a Média das IES privadas ( 42,3 ).
O desvio-padrão para as IES públicas foi de 13,8, maior do que para as IES privadas que foi de 13,7.
No tocante à organização acadêmica, os estudantes de centros universitários obtiveram a média de 42, enquanto os de universidades, faculdades e CEFET/IF obtiveram, respectivamente, Média de 42,7, 48,3 e 49,9.
A média da modalidade a distância ( 42,3 ) foi inferior à média da modalidade presencial ( 48,6 ). 7.1.2 ESTATÍSTICAS BÁSICAS NO COMPONENTE DE FORMAÇÃO GERAL Na Tabela 7.3, são apresentadas as Estatísticas básicas em relação ao componente da prova que avalia a Formação Geral dos estudantes concluintes.
A nota de Formação Geral é obtida levando -se em conta os acertos nas oito questões objetivas e a nota média obtida nas questões discursivas 1 e 2, sendo a nota de cada questão discursiva obtida levando -se em conta 80 \% da nota relativa ao conteúdo e 20 \% da nota relativa ao desempenho linguístico.
Conforme os dados da Tabela 7.3, os estudantes de todo o Brasil obtiveram, no Componente de Formação Geral da prova, desempenho médio de 45,3.
Quanto à variabilidade, o Desvio-padrão das notas dos estudantes do Brasil foi 17,2.
A maior Média foi obtida na região Sudeste ( 46,4 ) e a menor, na região Centro-Oeste ( 43,4 ).
As demais Médias foram : 44,3 na região Nordeste, 43,5 na região Norte e 44,4 na região Sul.
Já o maior Desvio-padrão foi obtido na região Nordeste ( 17,8 ), e o menor, na região Centro-Oeste ( 16 ).
Os demais desvios foram : 16,8 na região Norte, 17,1 na região Sul e 17,4 na região Sudeste.
A maior nota no Componente de Formação Geral da prova do Enade/2022 foi 98,2, obtida por pelo menos um estudante na região Nordeste.
A menor nota Máxima foi obtida na região Norte ( 81,2 ).
A Mediana do Brasil, como um todo, foi de 45,1 ; sendo a menor Mediana encontrada na região Centro-Oeste ( 42,9 ), e a maior, na região Sudeste ( 46,3 ).
A nota Mínima nessa parte foi zero em quatro das cinco regiões.
No Gráfico 7.2, é possível verificar a avaliação do desempenho dos estudantes do Componente de Formação Geral, a partir do histograma da distribuição das notas correspondentes.
Nesse gráfico, a distribuição é unimodal, com moda em ( 40 ; 50 ].
Na Tabela 7.4, são apresentadas as informações referentes ao desempenho dos concluintes do Componente de Formação Geral, em diferentes agregações : categoria administrativa, organização acadêmica e modalidade de oferta.
A partir dos dados da Tabela 7.4, observa -se que as Médias dos tipos de organização acadêmica foram : 44 para as universidades ; 43 para os centros universitários, 50,3 para as faculdades e 52,5 para os CEFET/IF.
No que se refere à modalidade de oferta, os cursos presenciais apresentaram Média superior ( 51 ) à Média dos cursos a distância ( 43,4 ).
A organização acadêmica que apresenta maior Máxima é o CEFET/IF com a marca de 98,2, as faculdades vêm logo em seguida com 96,3. 7.1.3 ESTATÍSTICAS BÁSICAS DO COMPONENTE DE CONHECIMENTO ESPECÍFICO Na Tabela 7.5, são apresentadas as estatísticas básicas referentes ao Componente de Conhecimento Específico da área de Tecnologia em Processos Gerenciais.
A nota do componente de Conhecimentos Específicos leva em conta as notas da parte objetiva e a média das notas obtidas nas três questões discursivas da parte discursiva da área.
De acordo com a Tabela 7.5, a média do desempenho dos estudantes do Brasil, para o Componente Específico da prova, foi 43,4.
A maior média foi obtida na região Sudeste ( 44,1 ), e a menor, na região Norte ( 40,3 ).
As demais Médias foram : 43,4 na região Sul, 42,6 na região Centro-Oeste e 40,9 na região Nordeste.
Quanto à variabilidade das notas, o desvio-padrão do Brasil foi 15,3 ; sendo o maior desvio-padrão observado na região Sudeste ( 15,6 ), e o menor, na região Norte ( 14,2 ).
Os demais desvios foram : 15,4 na região Nordeste, 15 na região Centro-oeste e 14,9 na região Sul.
A mediana das notas dos estudantes de todo o Brasil foi 43.
A maior mediana ocorreu na região Sudeste ( 44,4 ), e a menor, na região Norte ( 41,1 ).
As demais medianas foram : 42,9 na região Sul, 42,5 na região Centro-Oeste e 41,2 na região Nordeste.
A nota máxima do Brasil, como um todo, foi 95,6 ; sendo obtida por pelo menos um estudante na região Centro-Oeste.
As demais notas máximas foram : 90,9 na região Sudeste ; 85,1 na região Sul ; 82,9 na região Nordeste e 74,8 na região Norte.
A maior nota mínima foi de 3,6 na região Norte.
Assim como nos Gráficos 7.1 e 7.2, antes mostrados, no Gráfico 7.3, apresentado adiante, é possível ser feita uma avaliação do desempenho de concluintes de Tecnologia em Processos Gerenciais, em relação ao Componente de Conhecimento Específico, com um histograma da distribuição das notas correspondentes.
Esta também é uma distribuição unimodal, e o grupo modal é o ( 40 ; 50 ], o mesmo do grupo modal para a Formação Geral e para a Prova, de forma geral.
Na Tabela 7.6, é apresentada uma comparação dos resultados em relação a categoria administrativa, organização acadêmica e modalidade de oferta, agora, levando -se em conta o desempenho dos estudantes do Componente de Conhecimento Específico da prova.
De acordo com os dados expostos na Tabela 7.6, no que se refere à organização acadêmica, a maior Média foi a das CEFET/IF ( 49 ), vindo a seguir a das faculdades ( 47,6 ), seguida pela das universidades ( 42,3 ) e, por fim, pelos centros universitários ( 41,6 ).
O maior desvio-padrão foi o dos CEFET/IF ( 15,5 ).
As Universidades obtiveram a maior nota máxima ( 95,6 ).
Os centros universitários obtiveram nota máxima 85,1 ; as faculdades, 90,9 ; e os CEFET/IF, 85,1.
As Medianas foram : 42,5 nas universidades, 41,4 nos centros universitários, 48,5 nas faculdades e 49,6 nos CEFET/IF.
A maior nota mínima foi 3,6, nos CEFET/IF.
Quanto à categoria administrativa, observa -se que as IES privadas tiveram nota Máxima maior que as IES públicas, obtendo assim, respectivamente, 95,6 e 90,9. 7.2 ANÁLISE DAS QUESTÕES OBJETIVAS Esta seção apresenta estatísticas selecionadas e histogramas das questões objetivas de Formação Geral ( 7.2.1 ) e de Conhecimento Específico ( 7.2.2 ).
São, também, apresentadas e comparadas as médias das subpopulações caracterizadas por grande região.12 7.2.1 COMPONENTE DE FORMAÇÃO GERAL Na Tabela 7.7, são apresentadas as estatísticas básicas relativas às oito questões objetivas do Componente da prova que abrange a Formação Geral dos estudantes.
De acordo com os dados da Tabela 7.7, em relação ao Componente de Formação Geral, a Média do Brasil foi 45,8.
A menor Média foi encontrada na região Centro-oeste ( 43,3 ), e a maior, na região Sudeste ( 46,9 ).
As demais Médias foram : 46,1 na região Nordeste, 44,8 na região Sul e 43,9 na região Norte.
O Desvio-padrão do Brasil foi 20,5 ; sendo os maiores desvios padrão encontrados na região Nordeste, Sudeste e Sul ( 20,6 ), e o menor desvio-padrão encontrado na região Norte ( 19,1 ). 12 As fórmulas utilizadas para o coeficiente de correlação ponto-bisserial e para o índice de facilidade estão apresentadas no Capítulo 2.
As Medianas do Brasil, como um todo, e de três regiões ( Sul, Nordeste e Sudeste ) foram iguais a 50,0.
Nas regiões Centro-oeste e Norte, a mediana foi menor ( 37,5 ).
A nota Máxima ( 100,0 ) foi alcançada em quatro grandes regiões : Sul, Sudeste, Nordeste e Centro-oeste.
As notas Mínimas ( 0,0 ) foram iguais para todas as regiões.
Na Tabela 7.8, exposta mais adiante, são apresentados o Índice de Facilidade e o Índice de Discriminação ( Ponto-Bisserial ) para cada uma das questões objetivas do Componente de Formação Geral.
De acordo com os dados da Tabela 7.8, as questões objetivas do Componente de Formação Geral, segundo o Índice de Facilidade, foram assim avaliadas : nenhuma teve o Índice de Facilidade classificado como “ Muito fácil ”.
Duas questões foram classificadas como fáceis, com valor de 0,61 e 0,78.
Três questões foram classificadas como “ Médio ”, por terem valores de 0,44, 0,45 e 0,58.
Segundo o Índice de Discriminação, sete questões foram classificadas como “ Muito bom ”, com valores entre 0,40 e 0,51.
Uma questão foi classificada como Bom, com índice de acerto de 0,37.
O Índice de Facilidade variou de 0,21 a 0,78 e o de Discriminação, de 0,37 a 0,51.
Na Tabela 7.9, é apresentada a distribuição das questões, levando -se em conta, simultaneamente, a classificação dos dois índices.
Os dados da Tabela 7.9 atestam que 7 questões que tiveram Índice de Discriminação Muito bom figuraram em três níveis de dificuldade – Fácil, Médio e Difícil : duas classificadas na categoria “ Fácil ”, três classificadas na categoria “ Médio ” e duas, na categoria “ Difícil ”.
A questão que teve Índice de Discriminação “ Bom ” figurou no nível de dificuldade “ Difícil ”.
No Gráfico 7.4, cada uma das cinco curvas representa o percentual de respostas em determinada alternativa da questão, em função do número de acertos dos estudantes nessa parte da prova ( Formação Geral/Múltipla Escolha ), antes de possíveis eliminações pelo critério Ponto-Bisserial.
Os gráficos relativos às demais questões de Formação Geral constam do Anexo IV. 7.2.2 COMPONENTE ESPECÍFICO Na Tabela 7.10, são apresentadas as estatísticas básicas em relação às questões objetivas do Componente de Conhecimento Específico da prova de Tecnologia em Processos Gerenciais por grande região.
Como se verifica a partir dos dados da Tabela 7.10, a Média do Brasil para o Componente de Conhecimento Específico da Prova foi 46,8.
A menor Média foi observada na região Norte ( 43,8 ), e a maior, na região Sudeste ( 47,4 ).
O desvio-padrão de todo o Brasil foi 16,3, sendo o maior desvio-padrão encontrado nas regiões Nordeste ( 16,6 ), e o menor, na região Norte ( 15,3 ).
A Mediana de todo o Brasil foi 45,8, o mesmo valor da Mediana encontrada nas regiões Centro-oeste, Norte, Nordeste e Sul.
A nota Máxima da prova foi 95,8, obtida nas questões objetivas do Componente de Conhecimento Específico, por pelo menos um estudante da região Sudeste e Centro-oeste.
Na região Norte, a nota máxima foi 83,3.
A nota Mínima foi zero, obtida por pelo menos um estudante nas regiões Centro-oeste, Nordeste, Sudeste e Sul.
Na Tabela 7.11, são apresentados os Índices de Facilidade e de Discriminação ( Ponto-Bisserial ) das questões objetivas do Componente de Conhecimento Específico, para os estudantes de Tecnologia em Processos Gerenciais.
Para facilitar a diferenciação das questões, serão usadas as mesmas cores da Tabela 7.8, anteriormente apresentada, para as diferentes classificações dos Índices de Facilidade e de Discriminação.
A partir do Índice de Facilidade obtido, apresentado na Tabela 7.11, pode -se concluir que a maioria das questões objetivas da prova foram consideradas de nível Difícil ( 14 ) e 5 questões foram classificadas como “ Fácil ”.
Ao considerar o Índice de Discriminação, sete questões foram consideradas “ Bom ”, 4 Fraco e 7 “ Muito Bom ”.
O Índice de Facilidade variou de 0,17 a 0,73, e o de Discriminação, de 0,14 a 0,49.
Na Tabela 7.12, é apresentada a distribuição das questões levando -se em conta, simultaneamente, a classificação dos dois índices.
Os dados da Tabela 7.12 atestam que 7 questões que tiveram Índice de Discriminação Bom figuraram em dois níveis de dificuldade – Médio e Difícil : três classificadas na categoria Médio e quatro, na categoria Difícil.
Atestam, também, que aqueles que tiveram Índice de Discriminação Muito bom figuraram em três níveis de dificuldade – Fácil, Médio e Difícil : três classificadas na categoria Fácil, uma classificada na categoria Médio e três classificadas na categoria Difícil.
A título de exemplo das análises do comportamento das questões objetivas, no Gráfico 7.5, analisa -se a Questão 33 do Componente de Conhecimento Específico.
No Gráfico 7.5, cada uma das cinco curvas representa o percentual de respostas em determinada alternativa da questão, em função do número de acertos dos estudantes nessa parte da prova ( Componente Específico ), antes de possíveis eliminações pelo critério Ponto-Bisserial.
Os gráficos relativos às demais questões do Componente Específico constam no Anexo IV. 7.3 ANÁLISE DAS QUESTÕES DISCURSIVAS Esta seção apresenta estatísticas selecionadas e histogramas das questões discursivas de Formação Geral ( 7.3.1 ) e Conhecimento Específico ( 7.3.2 ).
Expõe e compara, também, as médias de subpopulações, caracterizadas por grande região.
Cumpre notar que nem todas as questões passam pelo crivo dos corretores.
As respostas depois de digitalizadas passam por um pré-processamento de identificação automático.
Algumas destas questões são definidas como um “ branco automático ” ( que não são enviadas para os professores para correção, a não ser que na segunda fase, a de inspeção visual, seja constatada uma resposta ).
O critério usado neste pré-processamento, é verificar para cada questão, a partir da quantidade de PIXELS, o que provavelmente, é uma resposta em branco.
A partir da separação destes documentos que foram considerados em branco, eles são enviados para uma verificação visual, por uma equipe de colaboradores, devidamente capacitados.
Deste modo, quaisquer erros, deste pré-processamento, são encontrados, e as questões preenchidas são enviadas para a correção dos professores.
Eventualmente, algumas respostas em branco não são detectadas nesse pré-processamento automático.
A presença de pequenos riscos, ou manchas, pode fazer o programa não identificar uma questão como um possível branco.
Estas são enviadas para a correção, mas são classificadas como “ branco ” pelos professores, quando da correção. 7.3.1 COMPONENTE DE FORMAÇÃO GERAL As análises dos resultados de desempenho dos estudantes de Tecnologia em Processos Gerenciais, nas duas questões discursivas relativas à Formação Geral, encontram -se na Tabela 7.13 e no Gráfico 7.6, cujos dados são analisados de forma sequenciada.
No caso da Tabela 7.13, analisa -se a média das notas de cada Questão Discursiva de Formação Geral que leva em conta as correções realizadas por duas bancas : uma que analisa e pontua aspectos relativos ao conteúdo das respostas, seguindo o padrão de respostas divulgado pelo Inep ( ver Anexo IX ) ; outra que avalia o desempenho linguístico expresso pela redação das respostas, seguindo um padrão de respostas específico da área de Língua Portuguesa.
Na Tabela 7.13, observa -se que a nota Média nesse conjunto de questões foi ligeiramente inferior à obtida nas objetivas.
Os estudantes de todo o Brasil obtiveram, em Formação Geral, Média 44,5 nas questões discursivas.
Pode -se notar, também, que o Desvio-padrão foi de 23,7 nas questões discursivas.
A maior Média foi obtida na região Sudeste ( 45,7 ), e a menor, na região Nordeste ( 41,5 ).
A Mediana de todo o Brasil, neste componente, foi 49,3 e a maior Mediana foi a da região Norte, de 50,7.
A nota Máxima ( 100 ) foi obtida na região Sul.
A nota Mínima ( 0,0 ) foi a mesma em todas as regiões do Brasil.
No Gráfico 7.6, está representada a distribuição das notas nas questões discursivas do Componente de Formação Geral.
Como mostra o Gráfico 7.6, a moda da distribuição ocorre no intervalo [ 50 ; 60 ], com frequência de aproximadamente 20 \% dos participantes, sendo que, no intervalo [ 0 ; 10 ], inclui -se, além da nota zero, a frequência de estudantes que deixaram esse tipo de questão em branco.
A distribuição é sempre decrescente, do primeiro intervalo, [ 0 ; 10 ] para o segundo intervalo [ 10 ; 20 ], crescente do segundo [ 10 ; 20 ] até o sexto intervalo [ 50 ; 60 ] e volta a ser decrescente, do sexto intervalo, [ 50 ; 60 ] ao décimo intervalo [ 90 ; 100 ].
Na sequência, os resultados verificados para cada uma das questões discursivas de Formação Geral serão apresentados, estabelecendo -se relações com os conteúdos abordados em cada uma delas.
Os comentários das Bancas de docentes corretores, a respeito do observado na correção das respostas dos estudantes, suas impressões e conclusões serão apresentados junto à análise de cada questão.
Cumpre esclarecer que, tendo em vista que as questões discursivas de Formação Geral são padronizadas, ou seja, constam de todas as provas, os comentários da Banca são os mesmos para todas as carreiras acadêmicas, sendo direcionados a todos os estudantes que participaram do Enade/2022.
A seguir, serão analisados o desempenho linguístico dos estudantes da Área de Tecnologia em Processos Gerenciais nas duas questões discursivas de Formação Geral do Enade/2022, apresentando -se os resultados obtidos e os comentários da banca de correção para cada questão. 7.3.1.1 ANÁLISE DE CONTEÚDO DA QUESTÃO DISCURSIVA 1 DO COMPONENTE DE FORMAÇÃO GERAL Os dados de Tecnologia em Processos Gerenciais obtidos a partir das respostas à Questão 1 do Componente de Formação Geral, encontram -se na Tabela 7.14 e no Gráfico 7.7, exposto após a apresentação dos dados da referida tabela.
Verifica -se, a partir dos dados da Tabela 7.14, que, na Questão 1 do Componente de Formação Geral – de desempenho abaixo ao da outra questão do mesmo componente –, os estudantes de todo o Brasil tiveram Média 43.
A maior Média para a Questão 1 foi obtida na região Sudeste ( 43,9 ), e a menor, na região Nordeste ( 40 ).
Quanto à variabilidade das notas, o Desvio-padrão de todo o Brasil foi 25,6.
O menor Desvio-padrão foi obtido na região Centro-Oeste ( 24,4 ), e o maior Desvio-padrão foi obtido pela região Norte ( 26,8 ).
A Mediana do Brasil, como um todo, foi 45,5, sendo igual no Sul.
A nota Máxima da Questão Discursiva 1 foi 100,0, obtida por pelo menos um estudante nas regiões Centro-oeste, Nordeste, Sul e Sudeste.
Na região Norte, a nota Máxima foi 95,5.
Já a nota Mínima para todas as regiões foi zero.
No Gráfico 7.7, apresenta -se a distribuição das notas na Questão Discursiva 1 do Componente de Formação Geral.
O Gráfico 7.7 mostra que mais de 13 \% dos participantes deixaram em branco a resposta à questão discursiva 1.
A moda dessa distribuição ocorre no intervalo ( 50 ; 60 ] com uma frequência maior que 15 \% dos participantes. 7.3.1.2 COMENTÁRIOS SOBRE A CORREÇÃO DE CONTEÚDO DAS RESPOSTAS À QUESTÃO DISCURSIVA 1 A Questão 1 abordou um tema atual, com dois itens para os estudantes responderem : a ) “ discuta por que os investimentos em educação, ciência, tecnologia e inovação são estratégicos e contribuem para o desenvolvimento de um país ” ; b ) “ explique como o fomento público ao desenvolvimento científico pode atender à justiça social e à equidade, em contextos como o da pandemia de Covid-19 ”.
O tema contempla como os investimentos em educação, ciência, tecnologia e inovação ( ECTI ) cooperam para um país se desenvolver e como o fomento público, vinculado ao desenvolvimento científico, pode atender à justiça social, à equidade e à população em momentos de pandemia.
Esse tema é atual, pois o mundo vivenciou recentemente uma pandemia que, somente no Brasil, deixou mais de 700 mil mortos, sendo exaustivamente comentada na televisão e nos meios digitais.
Logo, almeja -se que os respondentes justifiquem a necessidade de investimentos em ECTI com foco em questões como : investimentos público em P\&D por meio de universidades públicas ; pesquisas realizadas em outros países que possibilitaram o acesso mais célere às vacinas e serviram de modelo para o que foi desenvolvido no Brasil ; pesquisadores de universidade, tanto brasileiras quanto estrangeiras como protagonistas para orientar e informar a respeito da Covid-19 por meio dos canais de comunicação ; desenvolvimento de equipamentos de segurança, como máscara de nariz ou de rosto, respiradores, que foram implementados e viabilizados por pesquisa e extensão em universidades e escolas técnicas.
Esperava -se que os candidatos soubessem o básico a respeito do tema, pois o assunto é bastante atual.
Quanto ao item B, esperava -se que os candidatos abordassem questões como : o papel das universidades e dos centros de pesquisa para a justiça e a equidade ; o investimento público em pesquisa, tecnologia e educação para garantir acesso igualitário da população aos recursos de saúde, podendo fazer um contraponto com os investimentos privados ; a política de distribuição de vacinas gratuitas por meio do Sistema Único de Saúde, que possibilitou o acesso às vacinas no Brasil ; e as questões econômicas mundiais que determinaram uma distribuição desigual de recursos de saúde em escala global, com os países mais ricos retendo a maior parte desses recursos.
Atualmente, é comprovado que os países dominantes são aqueles que se desenvolvem e dominam as ferramentas tecnológicas.
As pesquisas mostram que os países que, mesmo em crise, investiram em ciência, tecnologia e inovação são os que conseguiram se reerguer mais rapidamente e sobressair -se aos demais países em crises pandêmicas.
Portanto, a questão foi elaborada de maneira oportuna e adequada ao universo do estudante universitário brasileiro.
No item A, o estudante poderia focar em diversos argumentos, entre eles : como ocorreu o investimento público em pesquisa e desenvolvimento por meio de universidades públicas e institutos de pesquisa ; como as pesquisas realizadas em outros países possibilitaram acesso mais rápido às vacinas e serviram de modelo para o que foi desenvolvido no Brasil ; como o contingenciamento do Fundo Nacional de Desenvolvimento Científico e Tecnológico ( FNDCT ) prejudicou e ainda prejudica o Brasil, entre diversos outros assuntos.
No item B, era esperado que o estudante explicasse como o fomento público ao desenvolvimento científico pode atender à justiça social e à equidade, em contextos como o da pandemia de Covid-19.
Além disso, deveria abordar a importância das universidades com a justiça social ; como os países que investem nessa área têm grandes vantagens em relação a outros países que não investem ; a distribuição de vacina atende à justiça social por possibilitar que toda a população seja vacinada, independentemente de ser rico ou pobre, por exemplo.
O espaço disponível para a resposta dos dois itens foi adequado e a maioria das respostas ocupou parte das linhas, não utilizando todo o espaço.
Algumas respostas não apresentaram a separação dos itens A e B.
A maioria dos estudantes respondeu de modo semelhante ao que constava do padrão de resposta.
Foi recorrente a exposição do problema sem desenvolvê -lo com clareza, apenas de forma superficial.
A abordagem do candidato sobre o tema tornou -se, em alguns casos, muito genérica.
As respostas consideradas boas atenderam ao que a banca esperava.
Abordaram a importância de investir em CT\&I para desenvolver o país, desdobrando os argumentos e citando como isso poderia ocorrer.
No geral, as notas estavam abaixo do esperado.
Ficou evidente que a maioria dos estudantes possui um conhecimento sobre os dois itens apresentados.
A Questão 1 visou avaliar o desenvolvimento da argumentação de um tema relevante para a sociedade brasileira vigente.
Além disso, houve problemas estruturais nos textos, pois a grande maioria dos estudantes tem uma noção do assunto, mas muitos possuem dificuldade de construir um texto com coesão e coerência, ou seja, não conseguem concatenar os argumentos para defender um ponto de vista.
No geral, em relação à Questão 1, identificou -se que a principal deficiência na formação dos estudantes está relacionada ao aprofundamento do conteúdo e à capacidade de expressar, com clareza e coesão, sua opinião a respeito de determinado assunto. 7.3.1.3 ANÁLISE DE CONTEÚDO DA QUESTÃO DISCURSIVA 2 DO COMPONENTE DE FORMAÇÃO GERAL A Tabela 7.15 apresenta dados estatísticos a respeito das notas de conteúdo da Questão Discursiva 2 do Componente de Formação Geral, organizados por grande região de funcionamento dos cursos, no país.
A Tabela 7.15 mostra que o desempenho médio dos estudantes, na Questão Discursiva 2 ( Média 46 ), foi acima do obtido na Questão Discursiva 1 ( Média 43 ).
A região Nordeste foi aquela cuja Média foi menor ( 43 ), e a maior Média foi na região Sudeste ( 47,4 ).
Quanto à variabilidade das notas, o Desvio-padrão de todo o Brasil foi 27,8 superior ao obtido na Questão Discursiva 1 ( 25,6 ).
O maior desvio nessa questão foi obtido na região Norte ( 29,7 ), enquanto o menor foi obtido na região Centro-Oeste ( 26,7 ).
A Mediana de todo o Brasil foi 53.
As regiões Norte, Sul e Centro-Oeste obtiveram Mediana de 51.
A nota Máxima da Questão Discursiva 2 foi 100,0, obtida por pelo menos 3 de todas as regiões, com exceção das regiões Norte e Nordeste.
Já a nota Mínima para todas as regiões foi zero.
No Gráfico 7.8, mostra -se a distribuição das notas na Questão Discursiva 2 do Componente de Formação Geral.
O Gráfico 7.8 mostra que quase 17 \% dos participantes deixaram em branco a resposta à questão discursiva 2.
A moda dessa distribuição, exceto as questões em branco, ocorre no intervalo [ 50 ; 60 ] com uma frequência maior que 15 \% dos participantes. 7.3.1.4 COMENTÁRIOS SOBRE A CORREÇÃO DE CONTEÚDO DAS RESPOSTAS À QUESTÃO DISCURSIVA 2 A Questão Discursiva 2 abordou um tema importante, com o objetivo de analisar conhecimentos gerais do aluno.
O tema contemplou o conhecimento geral do patriarcado estrutural ( dominação masculina ), o qual é recorrente nos noticiários, e um infortúnio na vida dos brasileiros, pois trata -se de uma questão de violência contra as mulheres.
É recorrente assistir ou ler notícias sobre o domínio masculino na sociedade brasileira.
Questões como a participação da mulher na política ; a paridade de salários entre profissionais de sexos diferentes ; a desigualdade social e o tratamento hostil que é dado às mulheres, em diversas situações, tanto profissionais quanto familiares, devem ser de conhecimento geral do estudante.
Não se objetiva tratar questões de ordem específica para a solução deste problema, mas, apenas, discorrer e argumentar sobre ele.
Portanto, a questão foi elaborada de maneira oportuna e adequada para o universo do estudante universitário brasileiro.
A Questão está dividida em dois itens : A e B.
No item A, esperava -se que o estudante abordasse o patriarcado – dominação masculina que gera a divisão sexual do trabalho, a desigualdade econômica e política, entre homens e mulheres, a sub-representação das mulheres na política e em cargos de poder.
Esperava -se que o estudante ressaltasse, ainda, questões culturais, como a educação sexista e a mídia, que coisificam a imagem das mulheres e naturalizam a desigualdade entre os sexos, de forma a considerar os homens como dominantes, enquanto as mulheres são historicamente tidas como submissas e incapazes.
No item B, era esperado que o estudante apresentasse a importância do papel do Estado na implementação de políticas públicas voltadas à proteção da mulher, tais como na área da saúde, na educação não sexista, na assistência social, na segurança pública – delegacias femininas, defensoria pública especializada, casa de abrigo, juizado especial etc., bem como ONGs que atuam no combate à violência contra a mulher e ao feminicídio.
O espaço disponível para a resposta dos dois itens foi adequado e grande parte das respostas ocupou parte das linhas, não contemplando o espaço completo.
Com relação à linguagem utilizada pelos estudantes, as respostas estavam dentro do esperado.
Alguns estudantes utilizavam termos do jargão do cotidiano.
No item B, o estudante deveria propor duas ações do Estado para o enfrentamento da violência contra as mulheres.
Os alunos deveriam apresentar propostas e desenvolvê -las.
Muitos estudantes responderam apenas por meio da apresentação de ações do Estado, mas sem desenvolvê -las ou argumentá -las.
Em geral, os alunos apenas expuseram ações – algumas vezes, uma, duas e, até mesmo, três ações –, mas sem fazer relações entre elas ou sem desenvolvê -las.
As respostas dos estudantes demonstraram que a maioria apresenta um conhecimento geral a respeito do patriarcado estrutural e da trágica desigualdade a qual impõe os homens sobre as mulheres, bem como de possíveis ações do Estado para a proteção delas.
A Questão visou avaliar o desenvolvimento da argumentação de um tema relevante para a sociedade brasileira vigente. 7.3.1.5 CONSIDERAÇÕES FINAIS SOBRE A CORREÇÃO DE CONTEÚDO DAS QUESTÕES DISCURSIVAS DO COMPONENTE DE FORMAÇÃO GERAL As questões discursivas da Área de Formação Geral abordaram temas atuais e proporcionaram a avaliação de competências importantes dentre as previstas pelas suas Diretrizes Curriculares.
A análise das respostas revelou, principalmente, dificuldade dos estudantes em relação à seleção de argumentos sólidos para discorrer sobre o tema – que pode ser fruto de desinteresse ou de desconhecimento mais profundo do contexto.
Esses fatos parecem ter sido os principais responsáveis pelo resultado aquém do esperado em termos de avaliação quantitativa.
Mesmo os estudantes que responderam em conformidade com o padrão demonstraram um universo de argumentação muito limitado, pobre em vocabulário, com o recurso frequente de procurar nos textos introdutórios os elementos para a construção de suas frases.
Encadear e relacionar argumentos foi uma deficiência percebida.
As razões eram apresentadas de forma confusa, sem fluência e sem constituição de unidade, por meio de termos genéricos e vagos. 7.3.1.6 ANÁLISE DE LÍNGUA PORTUGUESA DAS QUESTÕES DISCURSIVAS DO COMPONENTE DE FORMAÇÃO GERAL Os dados de Tecnologia em Processos Gerenciais obtidos a partir das respostas às questões discursivas do Componente de Formação Geral, no que tange à Língua Portuguesa, são apresentados na Tabela 7.16.
Os dados da Tabela 7.16 mostram que, em relação ao quesito Língua Portuguesa, os estudantes de todo o Brasil obtiveram Média 54,9.
A maior média em relação à Língua Portuguesa foi obtida na região Sudeste ( 55,7 ), e a menor, na região Nordeste ( 51,9 ).
Quanto à variabilidade das notas, o desvio-padrão de todo o Brasil foi 27,7.
O menor desvio-padrão foi obtido na região Centro-Oeste ( 26,4 ), e o maior desvio-padrão foi obtido na região Norte ( 30,7 ).
A mediana das notas de Língua Portuguesa foi 65 para o Brasil, como um todo, a mesma obtida em todas as regiões, exceto no Centro-Oeste e no Nordeste ( 57,5 ).
A nota máxima para todo o Brasil foi 100, com pelo menos um estudante tirando essa nota em todas as regiões.
A nota mínima foi zero em todas as regiões do país.
No Gráfico 7.9, mostra -se a distribuição das notas de Língua Portuguesa do Componente de Formação Geral.
O Gráfico 7.9 mostra que 15 \% dos participantes não pontuaram em língua portuguesa nas questões discursivas do componente de formação geral.
Excetuando -se a não pontuação, a moda dessa distribuição ocorre no intervalo ( 60 ; 70 ] com uma frequência maior que 15 \% dos participantes. 7.3.1.7 COMENTÁRIOS SOBRE A CORREÇÃO DE LÍNGUA PORTUGUESA DAS RESPOSTAS ÀS QUESTÕES DISCURSIVAS Os comentários feitos neste item a respeito da correção de Língua Portuguesa das respostas emitidas para as questões discursivas do Componente de Formação Geral estão agrupados em três eixos : avaliação de desempenho linguístico ; padrão de resposta utilizado na avaliação e análise das respostas dadas pelos estudantes. a ) A avaliação de desempenho linguístico De acordo com a Portaria n.º 298 do MEC/INEP, de 30 de junho de 2022, o Enade deve avaliar se o estudante desenvolveu, no processo de formação, dentre outras, competências para [ … ] V – Compreender as linguagens e suas respectivas variações ; VI – Ler, produzir e interpretar textos com clareza e com coerência ; VII – Identificar, analisar e interpretar diferentes representações verbais, não verbais, gráficas e numéricas de fenômenos diversos ou de um mesmo significado ; VIII – Formular e articular argumentos e contra-argumentos consistentes em situações sociocomunicativas ; IX – Identificar, compreender e analisar situações-problema utilizando pensamento holístico e sistêmico ao se abordar a complexidade da realidade. ( INEP, 2022d ) Essas competências são essenciais à formação de nível superior, porque dizem respeito ao domínio da linguagem, tanto no processo de leitura quanto na articulação e na organização de ideias para a produção textual.
Dessa forma, as duas questões discursivas de formação geral oferecem oportunidade para avaliar se o participante consegue desenvolver um tema, recorrendo a argumentos que explicitem sua opinião e elaborando propostas de intervenção social adequadas a um exercício consciente da cidadania.
É, portanto, necessário identificar, nas respostas discursivas de Formação Geral, as diferentes habilidades necessárias ao bom desempenho linguístico na situação de formalidade exigida : a organização das ideias, tendo em vista a construção de um texto coerente e estruturado ; o domínio das convenções ortográficas e das convenções morfossintáticas que constituem o arcabouço da gramática da norma-padrão da Língua Portuguesa ; e o domínio de um vocabulário amplo e diversificado, que permita uma comunicação clara e eficiente.
A Questão Discursiva 1 desenvolveu o tema de investimentos públicos em educação, ciência, tecnologia e inovação, que são estratégicos e contribuem para o desenvolvimento científico de um país, podendo abordar essa questão tanto no Brasil quanto em outros países.
O enunciado estava claro ao encaminhar a solicitação de resposta, evidenciando a exigência de elaboração de um texto que discutisse e explicasse por que os investimentos públicos em educação, ciência, tecnologia e inovação são estratégicos e contribuem para o desenvolvimento científico de um país e como o fomento público ao desenvolvimento científico pode atender à justiça social e à equidade, em contextos como o da pandemia de Covid-19.
Os três textos motivadores favoreceram a formulação e a articulação de argumentos e contra-argumentos consistentes, em defesa do ponto de vista a ser defendido pelo participante.
A Questão Discursiva 2 desenvolveu o conceito de patriarcado, tomando como base um texto de Biroli.
Com base nesse texto, o comando solicitou do participante :"discorra sobre a relação entre o patriarcado estrutural e a desigualdade entre homens e mulheres ” e “ proponha duas ações do Estado para o enfrentamento dos diferentes tipos de violência contra a mulher ”.
O comando “ discorra ” demandava a elaboração de um texto de base dissertativa, com desenvolvimento de argumentação sobre o tema.
Em relação à segunda parte, o comando exigia a elaboração de duas propostas de ações, que deveriam conter informações sobre o agente da intervenção ( Estado ) para o enfrentamento da situação apresentada.
As duas questões, portanto, assemelhavam -se, pois solicitavam um texto dissertativo ao apresentarem comandos para “ discutir ” e “ discorrer ”.
A Questão Discursiva 2, entretanto, solicitava que o aluno apresentasse propostas de ações para o enfrentamento do problema destacado.
Esses dois estilos textuais ofereceram oportunidade de avaliar se os participantes estavam habilitados a desenvolverem textos com esses dois distintos objetivos comunicativos, o que seria esperado ao final do curso de graduação e necessário à sua vida profissional.
A análise do desempenho linguístico esperado de um estudante que está concluindo seu curso superior apoiou -se no princípio de que ele deveria utilizar conhecimentos sobre os temas e estruturar seus textos de acordo com as características do registro formal adequado à situação comunicativa – avaliação de conhecimentos ao final de um curso superior.
Essa configuração determina exigências quanto à adequação da seleção vocabular, ao desenvolvimento do conteúdo, à estruturação sintática dos períodos, à organização lógica das ideias, à utilização de procedimentos de encadeamento textual e de referenciação, à obediência às exigências morfossintáticas próprias da modalidade escrita da norma-padrão, ao respeito às regras ortográficas de acentuação gráfica e de grafia das palavras. a ) O padrão de resposta utilizado na avaliação Em relação a este eixo de análise, foram considerados, no padrão de resposta, as seguintes competências : ( 1 ) estruturação textual ; ( 2 ) respeito às convenções ortográficas da norma-padrão da Língua Portuguesa ; ( 3 ) domínio dos diferentes aspectos morfossintáticos próprios da modalidade escrita formal da norma-padrão da Língua Portuguesa ; e ( 4 ) seleção vocabular adequada à modalidade escrita estratégicos e contribuem para o desenvolvimento científico de um país e como o fomento público ao desenvolvimento científico pode atender à justiça social e à equidade, em contextos como o da pandemia de Covid-19.
Os três textos motivadores favoreceram a formulação e a articulação de argumentos e contra-argumentos consistentes, em defesa do ponto de vista a ser defendido pelo participante.
A Questão Discursiva 2 desenvolveu o conceito de patriarcado, tomando como base um texto de Biroli.
Com base nesse texto, o comando solicitou do participante :"discorra sobre a relação entre o patriarcado estrutural e a desigualdade entre homens e mulheres ” e “ proponha duas ações do Estado para o enfrentamento dos diferentes tipos de violência contra a mulher ”.
O comando “ discorra ” demandava a elaboração de um texto de base dissertativa, com desenvolvimento de argumentação sobre o tema.
Em relação à segunda parte, o comando exigia a elaboração de duas propostas de ações, que deveriam conter informações sobre o agente da intervenção ( Estado ) para o enfrentamento da situação apresentada.
As duas questões, portanto, assemelhavam -se, pois solicitavam um texto dissertativo ao apresentarem comandos para “ discutir ” e “ discorrer ”.
A Questão Discursiva 2, entretanto, solicitava que o aluno apresentasse propostas de ações para o enfrentamento do problema destacado.
Esses dois estilos textuais ofereceram oportunidade de avaliar se os participantes estavam habilitados a desenvolverem textos com esses dois distintos objetivos comunicativos, o que seria esperado ao final do curso de graduação e necessário à sua vida profissional.
A análise do desempenho linguístico esperado de um estudante que está concluindo seu curso superior apoiou -se no princípio de que ele deveria utilizar conhecimentos sobre os temas e estruturar seus textos de acordo com as características do registro formal adequado à situação comunicativa – avaliação de conhecimentos ao final de um curso superior.
Essa configuração determina exigências quanto à adequação da seleção vocabular, ao desenvolvimento do conteúdo, à estruturação sintática dos períodos, à organização lógica das ideias, à utilização de procedimentos de encadeamento textual e de referenciação, à obediência às exigências morfossintáticas próprias da modalidade escrita da norma-padrão, ao respeito às regras ortográficas de acentuação gráfica e de grafia das palavras. a ) O padrão de resposta utilizado na avaliação Em relação a este eixo de análise, foram considerados, no padrão de resposta, as seguintes competências : ( 1 ) estruturação textual ; ( 2 ) respeito às convenções ortográficas da norma-padrão da Língua Portuguesa ; ( 3 ) domínio dos diferentes aspectos morfossintáticos próprios da modalidade escrita formal da norma-padrão da Língua Portuguesa ; e ( 4 ) seleção vocabular adequada à modalidade escrita formal da Língua Portuguesa, exigida pela situação comunicativa.
A seguir, são apresentados os aspectos avaliados associados a cada uma destas competências. ( 1 ) Estruturação textual condizente com o gênero textual solicitado e o modo de organização textual adequado ao gênero ( argumentativo ou expositivo ).
Esta competência envolve : - estruturação sintática condizente com o padrão da modalidade escrita formal da Língua Portuguesa, de modo a garantir a clareza e a coerência necessárias ; - distribuição do conteúdo do texto em parágrafos, de modo a garantir a sua organização temática ; - distribuição das ideias em períodos que mantenham uma estruturação sintática condizente com a modalidade formal da língua ( evitando -se frases fragmentadas, sequência justaposta de ideias sem encaixamentos sintáticos e presença de oração subordinada sem oração principal que reflitam hábitos da oralidade e comprometam a estrutura lógico-gramatical do texto ) ; - utilização de operadores discursivos que contribuam para a progressão temática do texto, estabelecendo relações lógicas entre as ideias apresentadas, tanto do ponto de vista intrafrasal como interfrasal : preposições, conjunções, pronomes relativos, advérbios e locuções adverbiais de tempo e lugar ; - utilização de procedimentos de referenciação lexical e pronominal que permitam a retomada de referentes textuais, como pronomes, advérbios e artigos ; além de elementos lexicais, como sinônimos, hiperônimos, hipônimos, expressões metafóricas e expressões resumitivas ; - emprego adequado dos pronomes relativos “ cujo(a) ” e “ onde ”, responsáveis por encaixamentos sintáticos no período composto ; - utilização adequada de sinais de pontuação que contribuam para a organização lógica da frase e do texto ; - inteligibilidade relacionada ao atendimento das exigências de estruturação textual.
A expectativa, portanto, era de que o participante recorresse a procedimentos linguístico-discursivos para organizar seu texto, permitindo o encadeamento lógico entre suas partes de forma a garantir a progressão e a coerência textuais, com estruturação que não reproduzisse hábitos da oralidade, e sim atendesse à norma-padrão da Língua Portuguesa. ( 2 ) Respeito às convenções ortográficas da norma-padrão da Língua Portuguesa.
Esta competência envolve o respeito às convenções estabelecidas pela legislação em vigor e consubstanciadas no Vocabulário Ortográfico da Língua Portuguesa, editado pela Academia Brasileira de Letras e avaliam -se : - domínio das regras de acentuação gráfica ; - domínio da grafia padrão das palavras ( com ausência de abreviaturas próprias da linguagem da internet ).
Assim, esperava -se que o participante : grafasse corretamente as palavras ; respeitasse as regras de acentuação gráfica ; empregasse maiúsculas em início de frase, em nomes próprios de pessoas, lugares ou instituições ; grafasse as siglas com letras maiúsculas ou apenas com a primeira letra maiúscula quando formarem uma palavra ; evitasse abreviações como p/, vc, tb, pq, tá, né, usadas muitas vezes em escrita informal e na internet ; obedecesse às regras de separação de sílabas no final da linha ; atendesse aos padrões de segmentação vocabular da norma-padrão. ( 3 ) Domínio dos diferentes aspectos morfossintáticos próprios da modalidade escrita formal da norma-padrão da Língua Portuguesa.
Esta competência envolve o domínio das convenções próprias da norma padrão da Língua Portuguesa relativas aos seguintes aspectos : concordância nominal ; concordância verbal ; regência nominal ; regência verbal ; flexão nominal ; flexão verbal ; correlação entre os tempos verbais ; colocação pronominal. ( 4 ) Assim, esperava -se que o participante : - flexionasse o verbo para estabelecer concordância de número com o sujeito da frase ; - flexionasse o artigo, o adjetivo e o pronome para concordar em número e em gênero com o substantivo a que cada um se refere ; - observasse a regência nominal e a verbal, utilizando a preposição adequada depois de um substantivo, um verbo ou um adjetivo ; - empregasse adequadamente o acento grave indicador da crase que expresse a contração entre uma preposição e um artigo, de acordo com as exigências de regência nominal ou verbal ; - obedecesse às regras de colocação pronominal ( próclise e ênclise ), distintas dos hábitos da oralidade ou da escrita informal ; - flexionasse adequadamente verbos, substantivos, adjetivos e pronomes, no que diz respeito à expressão das categorias gramaticais ; - flexionasse os verbos para expressar a correlação de modo e de tempo nas estruturas subordinadas.
Com base nesses critérios, alguns desvios ortográficos foram considerados como desvios de caráter morfossintático porque interferem na classificação dos vocábulos de base morfológica e funcional : - eliminação da marca de infinitivo ( -r- ) e substituição por acento agudo ou ausência total de marca do infinitivo, causando ambiguidade entre distintas classes de palavras ; - confusão entre as terminações “ ão ” e “ am ”, utilizadas para marcar diferentes tempos verbais ; - emprego indevido da forma verbal de"haver"\&lt ; “ há”\&gt ; no lugar da preposição ou do artigo \&lt ; “ a”\&gt ;, e vice-versa ; - emprego de hífen para separar pronome átono – tanto uso indevido quanto omissão ( ex : “ esperasse ”, em lugar de “ espera -se ” ; “ falar -mos ”, no lugar de “ falarmos ” ) ; - emprego de “ esta ” no lugar de “ está ”, provocando ambiguidade devido à confluência gráfica entre o pronome demonstrativo e a forma do verbo"estar", pela ausência da acentuação gráfica ; - emprego da conjunção aditiva"e"no lugar da forma verbal “ é ”, provocando ambiguidade devido à confluência gráfica por ausência da acentuação gráfica ; - emprego do advérbio “ mais ” no lugar da conjunção adversativa “ mas ”, reproduzindo confluência fônica da oralidade ; - registro da terceira pessoa do plural dos verbos “ ter ” e “ vir ” sem o acento circunflexo, causando desvio de concordância verbal ; - emprego inadequado do relativo “ onde ” deslocado do seu sentido locativo, reproduzindo tendência da oralidade ; - emprego do pronome relativo “ que ” sem a necessária preposição ( desvio de regência ) ; - ausência do acento grave obrigatório para marcar a crase. ( 5 ) Seleção vocabular adequada à modalidade escrita formal da Língua Portuguesa, exigida pela situação comunicativa.
Esta competência envolve : - precisão na seleção/utilização do vocabulário relacionado à temática solicitada pela questão ; - diversificação e adequação vocabular ; - coerência na combinação vocabular ; - ausência de marcas da oralidade, como termos de sentido muito genérico ( “ coisa ”, “ negócio ”, “ você ” ) e termos de registros mais informais ( como gírias, jargões, frases feitas, ditados populares, termos regionais ).
Assim, esperava -se que o participante : respeitasse a adequação vocabular, não empregando gírias ou expressões coloquiais ; apresentasse repertório vocabular diversificado e preciso ; evitasse repetição desnecessária de palavras ; utilizasse um vocabulário mais formal, como solicitado por um texto dissertativo.
Essas quatro competências foram reunidas em três itens de avaliação, atendendo aos aspectos explicitados anteriormente : I - Aspectos Ortográficos ( 30 \% ) : domínio das convenções ortográficas : grafia de vogais e consoantes, uso de maiúsculas e minúsculas, emprego do hífen e da acentuação gráfica, segmentação de palavras ; II - Aspectos Textuais ( 20 \% ) : domínio dos procedimentos de estruturação textual : organização interna dos períodos, emprego de conectores para a articulação lógica entre os períodos e entre os parágrafos, emprego de marcas de referenciação lexical e pronominal, utilização dos sinais de pontuação que contribuem para a organização lógica da frase ; III - Aspectos Morfossintáticos e Vocabulares ( 50 \% ) : domínio das regras de caráter morfossintático estabelecidas como modelares do ponto de vista da modalidade escrita formal da norma-padrão da Língua Portuguesa : concordância nominal e verbal, regência nominal e verbal, colocação pronominal, flexão nominal e verbal, correlação entre tempos e modos verbais, ausência de marcas de oralidade.
Outra competência combinada a esta na análise das respostas é a seleção vocabular adequada à modalidade escrita formal da Língua Portuguesa, tendo em vista a intersecção entre as duas do ponto de vista das exigências do registro formal da modalidade escrita da norma-padrão. b ) Análise das respostas dadas pelos estudantes A Questão 1 propiciou a oportunidade de maior desenvolvimento do tema solicitado, gerando consequências na estruturação textual das respostas, já que os textos foram mais expositivos, mais elaborados e fluentes, inclusive, um dos textos era uma tabela.
Já a Questão 2 teve um comportamento distinto, devido ao seu caráter mais opinativo, o texto foi mais curto e o vocabulário mais limitado.
Observaram -se, nas duas questões, alguns casos de comprometimento de autoria, com textos muito próximos a paráfrases dos textos motivadores e sem marcas de repertório sociocultural produtivo.
Apresentam -se, a seguir, algumas observações sobre o desempenho dos participantes : I – Aspectos Ortográficos : o desempenho dos participantes revelou uma diferença muito grande nos dois aspectos analisados, como já foi observado nos relatórios dos anos anteriores : baixo índice de desvios da grafia padrão e grande índice de desvios de acentuação.
Em vários casos, ocorreu ausência quase completa de acentuação gráfica.
Vale ressaltar que muitos participantes apresentaram uma caligrafia quase ininteligível.
Ainda nesse aspecto, ocorreu a omissão inadequada de cedilha, de til e do traço do T em algumas respostas.
Houve, inclusive, respostas com ausência total de cedilha ou til.
Os casos mais sistemáticos de eliminação do acento indicador da sílaba tônica foram : palavras proparoxítonas ( “ politica ”, “ pratica ”, “ publica ”,"policia") ; palavras paroxítonas terminadas em ditongo crescente ( “ necessario ”,"historia","propria", “ violencia ” ) ; palavras oxítonas ( “ tambem ”,"populacao", “ pais ” ).
Por outro lado, destaca -se o uso indevido do acento gráfico em palavras como “ investímento ” e “ pátriarcado ”, por exemplo.
Quanto aos desvios de grafia, observam -se os seguintes exemplos : emprego inadequado de maiúscula : “ Leis ” ; ou de minúscula :"constituição"“ Ministério publico ” ; confusão entre"ss","sc","c"e"ç": “ nessessaria", “ agreção ”, “ conciencia ”, “ coMECam ”, “ cocientização ”, “ pacos ” ; eliminação de grafemas : “ mostando ”,"deveriamo"; - desvios de segmentação : “ des de", “ des igualdade ”, “ pre conceito"; troca de grafema : “ órgam ” ( por “ órgão ” ), “ escravido ” ( por “ escravidão ” ).
II – Aspectos Textuais : a estruturação textual foi a competência que apresentou o maior número de desvios, sendo responsável, em alguns textos, pela falta de clareza e pela desorganização.
No conjunto dos textos analisados, observaram -se diferentes níveis de domínio dos recursos de estruturação e organização textual.
No entanto, os participantes podem ser agrupados em três faixas descritas a seguir : - os que possuem grande domínio dos recursos de estruturação textual, com produção de textos que atendem às exigências de clareza e coerência, com progressão temática bem direcionada e utilização produtiva de operadores argumentativos ; - os que apresentam um domínio regular dos recursos de estruturação textual, com desvios na organização do texto e ausência de paragrafação, além de uso insuficiente de operadores argumentativos e processos de encaixamento sintático ; - os que apresentam baixo domínio dos recursos de estruturação textual, com muitos desvios que retomam hábitos da oralidade e que comprometem a coerência do texto : ocorrência de sequência justaposta de ideias sem encaixamentos sintáticos ; redução de estruturas subordinadas, ao lado do aumento na frequência de estruturas coordenadas e absolutas ; redução no uso de conectores para expressar relações lógicas essenciais à construção do texto ; emprego equivocado de operadores que não estabelecem relações lógicas coerentes entre ideias do texto ; fragmentação frasal que compromete a estrutura lógico-gramatical.
Uma estrutura que vem se ampliando na escrita dos participantes, refletindo hábitos da oralidade, é a duplicação de um dos elementos da frase, ou seja, a retomada anafórica de um elemento topicalizado, como por exemplo :"Elas precisam saber o verdadeiro papel DELAS na sociedade “, A desigualdade ELA vem de uma cultura machista... ” ; “ Homem e mulher, ELES são iguais ”.
Também em relação ao emprego dos sinais de pontuação, observaram -se distintos tipos de desempenho : participantes que dominam a funcionalidade da pontuação para a construção de textos fluentes e coesos, por um lado, e um grande número de participantes que apresentam precariedade em sua utilização, chegando ao ponto de elaborar parágrafos sem qualquer marca interna de pontuação para separar os períodos.
São os seguintes os tipos de desvios encontrados : - emprego de vírgula para separar o sujeito e o predicado : “ A mulher, sofre violência ” ( no lugar de “ A mulher sofre violência ” ) ;"O correto seria, ter um canal de denúncia ” ( no lugar de “ o correto seria ter um canal de denúncia ” ) ; - ocorrência de apenas uma das vírgulas para separar uma palavra, uma expressão ou uma oração encaixada : “ O homem, por sua vez não aceita que uma mulher seja sua chefa ” ( no lugar de “ O homem, por sua vez, não aceita que uma mulher seja sua chefa ” ) ; - emprego de ponto no lugar de vírgula para separar orações de um mesmo período : “ A todo custo as mulheres são atacadas na política.
Sexualmente também. ” ( no lugar de “ A todo custo as mulheres são atacadas na política, sexualmente também. ” ) ; - ausência de vírgula ou conjunção para separar elementos de uma enumeração :"... conquistas como voto equidade salarial direito de estudar e reconhecimento..."( no lugar de"... conquistas como voto, equidade salarial, direito de estudar e reconhecimento...") ; - ausência de vírgula para separar oração adjetiva explicativa ou utilização inadequada para separar oração adjetiva restritiva : “ As mulheres, que são iguais aos homens devem ser igualmente protegidas por lei."( no lugar de “ As mulheres, que são iguais aos homens, devem ser igualmente protegidas por lei.") ; - emprego indevido do sinal de dois pontos : “ Mulheres são julgadas : pelas simples roupas que usam. ” ( no lugar de “ Mulheres são julgadas pelas simples roupas que usam. ” ) ; - ausência de pronome no processo de referenciação : “ julga a mulher ou até agride ” ( no lugar de “ julga a mulher ou até a agride ” ” ) ; - emprego de vírgula em lugar do ponto-final : “ O patriarcado estrutural é formado por elementos históricos e religiosos, há pouquíssimo tempo, as mulheres passaram a votar. ” ( no lugar de “ O patriarcado estrutural é formado por elementos históricos e religiosos.
Há pouquíssimo tempo, as mulheres passaram a votar. ” ).
III – Aspectos Morfossintáticos e Vocabulares : o domínio das convenções morfossintáticas apresentou resultados que evidenciam, claramente, alguns aspectos problemáticos no desempenho dos participantes.
Assim como na análise dos desvios textuais, pode -se observar uma distinção entre três grupos de participantes : o grupo que revela domínio das convenções morfossintáticas ; o grupo que apresenta um domínio regular, com desvios pontuais ; e o grupo que revela domínio precário dessas convenções.
Pode -se destacar a regência, tanto verbal quanto nominal, como o aspecto com maior número de desvios : a ) Regência nominal e verbal : Na maioria dos textos, destacam -se os desvios relacionados à noção de"regência", ou seja, à relação de subordinação que se estabelece entre as palavras por meio, ou não, de preposições, refletindo uma tendência na modalidade escrita da língua.
O desvio mais frequente foi a falta do sinal indicativo da crase, em contextos em que seria obrigatória, ou, ao contrário, um uso indevido em contextos em que o princípio não se aplica.
Os dois tipos de desvios revelam que o usuário não tem consciência de que, sob a forma do termo “ a ”, se pressupõe a presença de uma preposição “ a ”, exigida pela regência do termo anterior.
Exemplos : “ diferentes tipos de agressão a mulher ” ( no lugar de “ diferentes tipos de agressão à mulher ” ) ; “ Desenvolverem práticas econômicas e sociais as novas gerações"( no lugar de “ Desenvolverem práticas econômicas e sociais às novas gerações") ; “ Além de proteger à mulher... ” ( no lugar de “ Além de proteger a mulher... ” ).
Outro desvio relacionado à regência verbal e à nominal, encontrado frequentemente nas questões, é a ausência de preposição antes de pronome relativo, refletindo um processo generalizado na modalidade oral da língua, em situações de registro informal.
Exemplo :"outra forma é oferecer apoio psicológico quem sofreu qualquer tipo de abuso..."( no lugar de :"outra forma é oferecer apoio psicológico a quem sofreu qualquer tipo de abuso...") b ) Concordância nominal e verbal : Os desvios relativos à aplicação das regras de concordância verbal e nominal correspondem a fatos que têm sido amplamente estudados na modalidade oral da Língua Portuguesa.
Quanto à concordância de número, observou -se ausência de marca ( com sujeito anteposto ou posposto ), como por exemplo a ausência de acento circunflexo na forma plural do presente do indicativo do verbo “ ter ”, considerada como um desvio na concordância verbal e não na acentuação gráfica.
Exemplo :"as mulheres ainda tem como obrigação"( no lugar de"as mulheres têm como obrigação") ; “ as leis em vigor tem que ser cumpridas ” ( no lugar de “ as leis em vigor têm que ser cumpridas ” ).
Outro desvio na aplicação da regra da concordância de número foi o aparecimento da marca de plural em verbos ou adjetivos relacionados a núcleos substantivos no singular, devido à presença de modificadores no plural, evidenciando um processo de hipercorreção.
Esse processo tem se ampliado nas situações de oralidade, caracterizando uma tendência de variação da Língua Portuguesa.
Exemplos :"Mudança nas leis que aumentem a pena do agressor"( no lugar de"Mudança nas leis que aumente a pena do agressor") ;"o aumento dos feminicídios devem ser combatidos ” ( no lugar de"o aumento dos feminicídios deve ser combatido").
Quanto à concordância de gênero, vários casos foram observados, normalmente no âmbito de sintagmas nominais longos, em que o adjetivo está afastado do substantivo, conforme apontam as pesquisas sobre a modalidade oral. tipos de desvios revelam que o usuário não tem consciência de que, sob a forma do termo “ a ”, se pressupõe a presença de uma preposição “ a ”, exigida pela regência do termo anterior.
Exemplos : “ diferentes tipos de agressão a mulher ” ( no lugar de “ diferentes tipos de agressão à mulher ” ) ; “ Desenvolverem práticas econômicas e sociais as novas gerações"( no lugar de “ Desenvolverem práticas econômicas e sociais às novas gerações") ; “ Além de proteger à mulher... ” ( no lugar de “ Além de proteger a mulher... ” ).
Outro desvio relacionado à regência verbal e à nominal, encontrado frequentemente nas questões, é a ausência de preposição antes de pronome relativo, refletindo um processo generalizado na modalidade oral da língua, em situações de registro informal.
Exemplo :"outra forma é oferecer apoio psicológico quem sofreu qualquer tipo de abuso..."( no lugar de :"outra forma é oferecer apoio psicológico a quem sofreu qualquer tipo de abuso...") b ) Concordância nominal e verbal : Os desvios relativos à aplicação das regras de concordância verbal e nominal correspondem a fatos que têm sido amplamente estudados na modalidade oral da Língua Portuguesa.
Quanto à concordância de número, observou -se ausência de marca ( com sujeito anteposto ou posposto ), como por exemplo a ausência de acento circunflexo na forma plural do presente do indicativo do verbo “ ter ”, considerada como um desvio na concordância verbal e não na acentuação gráfica.
Exemplo :"as mulheres ainda tem como obrigação"( no lugar de"as mulheres têm como obrigação") ; “ as leis em vigor tem que ser cumpridas ” ( no lugar de “ as leis em vigor têm que ser cumpridas ” ).
Outro desvio na aplicação da regra da concordância de número foi o aparecimento da marca de plural em verbos ou adjetivos relacionados a núcleos substantivos no singular, devido à presença de modificadores no plural, evidenciando um processo de hipercorreção.
Esse processo tem se ampliado nas situações de oralidade, caracterizando uma tendência de variação da Língua Portuguesa.
Exemplos :"Mudança nas leis que aumentem a pena do agressor"( no lugar de"Mudança nas leis que aumente a pena do agressor") ;"o aumento dos feminicídios devem ser combatidos ” ( no lugar de"o aumento dos feminicídios deve ser combatido ” ).
Quanto à concordância de gênero, vários casos foram observados, normalmente no âmbito de sintagmas nominais longos, em que o adjetivo está afastado do substantivo, conforme apontam as pesquisas sobre a modalidade oral.
Exemplo : “ E as leis sobre as mulheres que sofrem violência ser mais rigorosos ” ( no lugar de “ E as leis sobre as mulheres que sofrem violência ser mais rigorosas ” ) ; Registro da forma verbal “ é ” sem acento, provocando ambiguidade na leitura pela confluência com a conjunção aditiva"e".
Exemplo : “ A desvantagem entre salários, poderes e entre a violência física e moral e crescente por demais... ” ( no lugar de “ A desvantagem entre salários, poderes e entre a violência física e moral é crescente por demais ”. c ) Emprego do pronome relativo “ onde ” como relativo universal, típica marca da oralidade, em que esse pronome não é mais utilizado apenas em sua função locativa"lugar onde se realiza a ação".
Exemplo :"Existe a desigualdade de gênero, onde homem ganha mais do que as mulheres"; “ há falta de informação e conhecimento onde ” ; “ O patriarcado está estruturado em nossa sociedade há séculos, onde o sexo masculino... ” d ) Falta de artigo definido antes de substantivo, comprometendo o paralelismo em enumerações : Exemplo :"seja em oportunidade de emprego, salário"( no lugar de"seja em oportunidade de emprego, de salário"). e ) Desvios relacionados à flexão verbal : Exemplos : emprego indevido do tempo ou do modo verbal : “ promovido ” por “ promovendo ” ; eliminação de marca de infinitivo : “ colocar ” por “ coloca ”, “ ta ” por “ estar ”, “ está ” por “ estar ”, “ esta ” por “ estar ” ; “ estás ” por “ estas ” ; utilização de construção de “ estar ” + gerúndio como marca de oralidade : “ a sociedade cada vez mais tá evoluindo ” ; segmentação indevida da desinência de número e pessoa de verbos : “ denunciar -mos ”.
Assim como nas demais competências avaliadas, há diferenças entre os participantes em relação à seleção e ao emprego do vocabulário.
Foram encontrados textos com alto padrão de domínio vocabular, chegando até ao emprego de vocabulário técnico das áreas do saber abordadas nas questões.
Por outro lado, na maioria dos textos, foram observados vários tipos de inadequação na seleção e na utilização do vocabulário, sendo que o principal deles foi a excessiva repetição de certas palavras, revelando limitação de repertório vocabular.
Além disso, observou -se a ocorrência de expressões da oralidade, revelando falta de domínio de vocabulário mais abstrato e de maior complexidade, essencial ao desenvolvimento do texto de base dissertativa.
Exemplos : seleção vocabular indevida : “ Em pleno século 21 é de emplena importância... ” ; emprego inadequado do verbo “ ter ” no lugar de “ haver ” : “ Durante a história, teve muito machismo ” ; imprecisão vocabular : “ Mas a cabeça, o corpo, há mesma igualdade para desenvolvermos igualmente as coisas ”, “ Em alguns anos atrás, as coisas eram mais complicadas nas questões de desigualdade no trabalho. ” A imprecisão vocabular, como o uso da palavra “ coisa ” em vários textos, demonstra o domínio precário da norma-padrão da Língua Portuguesa, com presença constante de traços de oralidade. 7.3.1.8 CONSIDERAÇÕES FINAIS SOBRE A CORREÇÃO DE LÍNGUA PORTUGUESA DAS RESPOSTAS ÀS QUESTÕES DISCURSIVAS DO COMPONENTE DE FORMAÇÃO GERAL O conjunto de textos analisados no Enade 2022 permite traçar um perfil bem detalhado do desempenho linguístico do estudante brasileiro de graduação.
Conforme ficou comprovado nos comentários dos exemplos de respostas apresentados neste relatório, uma parcela significativa dos participantes revela lacunas importantes em sua formação, principalmente no que diz respeito à estruturação e à organização de textos.
Como esta é uma competência essencial à produção escrita, que deve ser construída ao longo dos anos de escolaridade ( do ensino fundamental à educação ), sua recuperação é trabalhosa e depende de um esforço intenso de exposição a textos e de reflexão sobre os requisitos que garantem a clareza e o encadeamento das ideias.
A competência relacionada ao domínio do vocabulário, que deve ser diversificado e apropriado ao padrão formal da língua, também apresentou um desempenho insatisfatório pela maioria dos participantes, que se limitam ao senso comum, com vocabulário genérico e vago, o que é inadequado para atender ao perfil de um graduando.
As outras duas competências avaliadas ( convenções ortográficas, convenções morfossintáticas ) apresentam desvios pontuais e são de mais fácil recuperação, com exceção dos desvios morfossintáticos que já estão incorporados na modalidade oral informal da Língua Portuguesa, configurando processos de variação linguística e, talvez, antecipando futuras mudanças no padrão da modalidade escrita.
Com base na avaliação do desempenho linguístico aqui relatada, conclui -se que é preciso aperfeiçoar a formação superior, oportunizando aos estudantes mais acesso às ferramentas necessárias para que sejam bem-sucedidos nas diferentes situações de sua vida cidadã – profissional, cultural e política.
Essa é uma contribuição efetiva que os resultados do Enade podem oferecer às IES do país.
Para finalizar, vale observar que as duas questões discursivas atenderam às exigências de avaliação, tanto do domínio discursivo dos estudantes quanto do seu conhecimento sobre temas relacionados ao conhecimento acadêmico, o que pode ser extremamente útil para que os cursos superiores possam contribuir na formação de seus graduandos. 7.3.2 COMPONENTE ESPECÍFICO A Tabela 7.17 apresenta dados sobre questões discursivas do Componente de Conhecimento Específico.
A partir dos dados registrados na Tabela 17.7, observa -se que a Média foi mais baixa do que para as questões discursivas do Componente de Formação Geral.
Enquanto no Componente de Formação Geral, a Média para estudantes de Tecnologia em Processos Gerenciais de todo o Brasil foi 44,5 ; na parte de Componente Específico, a Média foi 23,9.
A maior média deste componente foi obtida pelos estudantes da região Sudeste ( 25,1 ), e a menor, pelos da região Nordeste ( 19,9 ).
Quanto à variabilidade das notas, o desvio-padrão de todo o Brasil foi 20,9.
O maior desvio-padrão foi encontrado na região Sudeste ( 21,4 ), e o menor, na região Norte ( 19,4 ).
No Gráfico 7.10, representa -se a distribuição das notas nas questões discursivas do Componente Específico.
À luz dos dados do Gráfico 7.10, verifica -se que a moda da distribuição das notas nas questões discursivas do Componente Específico ocorre no intervalo ( 20 ; 30 ], com mais que 10 \% do total de participantes, excetuando -se a não pontuação dos participantes.
A análise de cada uma destas questões será feita a seguir. 7.3.2.1 ANÁLISE DA QUESTÃO DISCURSIVA 3 DO COMPONENTE ESPECÍFICO Os resultados da Questão Discursiva 3 do Componente Específico são descritos na Tabela 7.18 e comentados na sequência.
Como mostrado na Tabela 7.18, na Questão 3, a média dos estudantes de todo o Brasil foi 17,4.
O desempenho dos estudantes nessa questão foi inferior ao das Questões 2 e 5.
A maior Média na questão em evidência foi obtida pelos estudantes da região Sudeste ( 18,9 ), enquanto a menor Média foi obtida na região Nordeste ( 13,6 ).
Quanto à variabilidade das notas, o desvio-padrão de todo o Brasil foi 21,8.
O maior desvio-padrão foi obtido na região Sudeste ( 22,7 ), enquanto o menor foi obtido na região Nordeste ( 19,8 ).
A nota máxima, 100,0 pontos, foi alcançada por pelo menos um estudante nas regiões Centro-Oeste, Sudeste e Sul.
A mediana do Brasil foi 10, para as regiões Centro-Oeste, Sudeste e Sul ; nas demais regiões foi 0.
A nota mínima zero foi a mesma em todas as regiões do Brasil.
Destaca -se que, com base nos dados do Gráfico 7.11, mais de 40,0 \% dos estudantes deixaram a resposta à questão em branco.
A moda dessa distribuição ocorre no intervalo ( 0 ; 10 ], com mais de 10 \% do total de participantes, excetuando -se as não pontuações. 7.3.2.2 COMENTÁRIOS SOBRE AS RESPOSTAS À QUESTÃO DISCURSIVA 3 A questão abordou dois temas importantes para a Gestão de Processos Gerenciais voltados para a indústria moveleira : o ciclo PDCA e a definição de um processo.
Tais temas podem ser ferramentas valiosas para qualquer tipo de indústria, permitindo a otimização dos processos de produção, a melhoria da qualidade dos produtos e o aumento da eficiência da cadeia de suprimentos, garantindo a satisfação do cliente final.
O tema faz parte da formação básica do curso e pertence às suas Diretrizes Curriculares Nacionais.
Esperava -se do estudante um conhecimento básico de como os temas abordados podem ser aplicados na gestão empresarial.
Com relação ao questionamento do primeiro item, era esperado que o estudando demonstrasse conhecimentos mínimos ( apesar de um pouco mais específicos ) a respeito das informações prévias sobre a metodologia de gestão.
Eram esperadas respostas mais diretas e objetivas, pois o item basicamente pedia para descrever o ciclo e demonstrar como ele poderia ser utilizado.
No que tange ao processo, segundo item, foi solicitado que o estudante mostrasse um conhecimento básico sobre a área, explicando -o e exemplificando suas atividades de entrada, processamento e saída.
Observou -se que os estudantes, de modo geral, não conseguiram reproduzir um conhecimento mais detalhado sobre o tema, em especial, sobre o ciclo PDCA e como ele pode ser utilizado na gestão de processos.
Em relação às etapas do modelo, uma pequena parte explanou essa resposta, mas sem demonstrar grande segurança e riqueza de detalhes.
Em relação à explicação clara e condizente de processo, poucos estudantes conseguiram expor da forma como o enunciado pedia, porém foi possível identificar um número maior de acerto, apesar das respostas curtas e diretas quanto às etapas do processo na indústria moveleira.
Com relação à linguagem utilizada pelos estudantes, as respostas estavam dentro do esperado.
Alguns estudantes utilizaram os termos técnicos de maneira muito adequada, enquanto outros usaram uma linguagem comum para exemplificar o que o enunciado pediu.
Tratando -se do item A, relativo ao ciclo PDCA e sua aplicação, notamos que as respostas preenchidas foram bem curtas e diretas, no que tange às etapas, pois os estudantes, na grande maioria, apenas citaram e deram um breve resumo de como a metodologia pode ser usada para gerir os processos.
Já sobre o item B, a grande maioria não soube explicar, de forma clara e objetiva, sobre o que se trata um processo, mas demonstrou as etapas : entrada, processamento e saída.
Os estudantes Exemplificaram, de forma bem resumida, demonstrando pouco conhecimento da Língua Portuguesa e suas regras gramaticais.
As respostas aos dois itens apresentaram uma pequena porção de informações que não estavam relacionadas ao contexto da questão, que pedia a demonstração de uma associação com a indústria moveleira.
Em alguns casos, foram citados exemplos, como a produção de bolos, a venda de frutas, entre outros assuntos não pertinentes ao tema em questão.
Portanto, houve falta de coerência entre o enunciado da questão e as respostas encontradas.
As respostas demonstraram que os estudantes tinham a noção básica sobre o item A no que tange a suas etapas e como podem ser utilizadas.
Além disso, sobre as etapas, na maioria, parece não ter havido dificuldade em serem citadas.
A maior dificuldade foi encontrada no item B, que exigia a explicação clara e condizente de processo.
As respostas consideradas fracas foram aquelas em que todos os itens foram respondidos erroneamente ou de forma parcial.
No caso do item A, as respostas medianas abordaram as etapas do ciclo PDCA de forma superficial, mencionando apenas algumas delas.
Já em relação ao item B, os estudantes não conseguiram explicar claramente o que é um processo.
As respostas consideradas boas foram aquelas em que os estudantes responderam completamente ou quase completamente aos itens solicitados na questão.
Ficou evidente, no item B, que os estudantes tinham o conhecimento das etapas de um processo para a fabricação de um móvel.
Porém, houve dificuldade em descrever, de modo objetivo, claro e condizente, o que é um processo.
Relativo ao ciclo PDCA, muitos sabiam as etapas, contudo, de forma abstrata e resumida. 7.3.2.3 ANÁLISE DA QUESTÃO DISCURSIVA 4 DO COMPONENTE DE CONHECIMENTO ESPECÍFICO Na Tabela 7.19, constam as informações relativas à Questão 4 do conjunto de Questões do Componente de Conhecimento Específico.
Como se observa a partir dos dados da Tabela 7.19, a Média geral do Brasil foi 24,2, sendo a maior Média registrada na região Sudeste ( 25,3 ), e a menor, na região Nordeste ( 20,7 ).
Quanto à variabilidade das notas, o desvio-padrão geral do Brasil foi 25,9 ; sendo o maior desvio-padrão registrado na região Sudeste ( 26,3 ), e o menor, na região Norte ( 25 ).
A nota máxima ( 100,0 ) foi atingida por pelo menos um concluinte nas regiões Centro-Oeste, Nordeste, Sudeste e Sul.
A Mediana em todo o Brasil foi 25 ; o mesmo valor foi obtido nas regiões Sudeste e Sul.
Nas regiões Centro-Oeste, Norte e Nordeste a Mediana foi de 0.
No Gráfico 7.12, está representada a distribuição de notas da Questão Discursiva 4, do Componente de Conhecimento Específico.
A distribuição mostrada no Gráfico 7.12, excetuando -se resultados nulos, tem moda no intervalo ( 40 ; 50 ] com mais que 25 \% do total de participantes.
Chama a atenção que mais de 45 \% dos participantes deixaram essa questão em branco. 7.3.2.4 COMENTÁRIOS SOBRE AS RESPOSTAS À QUESTÃO DISCURSIVA 4 A questão abordou dois temas importantes para a Gestão de Processos voltados para os fundamentos da Administração e o processo comunicacional.
Ambas as habilidades são fundamentais para garantir a eficiência e a eficácia das atividades de uma organização, promover um ambiente colaborativo e produtivo e alcançar os objetivos e as metas estabelecidas.
O tema faz parte da formação básica do curso e pertence às suas Diretrizes Curriculares Nacionais.
Esperava -se do estudante um conhecimento básico de como os temas abordados podem ser aplicados em uma organização.
Com relação ao questionamento do primeiro item, esperava -se que o estudante demonstrasse conhecimentos mínimos ( apesar de um pouco mais específico ) a respeito das informações prévias sobre a gestão empresarial.
Eram esperadas respostas mais diretas e objetivas, pois o item basicamente pedia para descrever cada modelo.
No que tange ao processo comunicacional, segundo item, o quesito solicitava que o estudante mostrasse um conhecimento básico sobre a área, explicando o processo e exemplificando -o com três elementos.
A maioria dos estudantes respondeu de forma semelhante ao que constava no padrão de resposta, pois não conseguiu reproduzir um conhecimento mais detalhado sobre o tema, em especial, sobre os elementos do processo comunicacional.
Com relação ao modo que como este modelo ocorre na administração participativa, uma pequena parte explanou a resposta, mas sem demonstrar grande segurança e riqueza de detalhes.
Já em relação à explicação clara e condizente dos fundamentos do modelo diretivo e participativo, poucos estudantes conseguiram se expressar conforme o enunciado pedia, porém foi possível identificar um número maior de acerto, apesar das respostas curtas e diretas quanto às definições básicas dos modelos.
Com relação à linguagem utilizada pelos estudantes, as respostas estavam dentro do esperado.
Alguns estudantes utilizaram os termos técnicos de maneira muito adequada, enquanto outros usaram uma linguagem comum para exemplificar o que o enunciado pedia.
Ao analisarmos o item A, que se refere ao modelo diretivo e participativo, percebemos que as respostas fornecidas foram bastante resumidas.
A maioria dos estudantes limitou -se a mencionar apenas algumas partes fundamentais de cada modelo ou apenas descreveu bem um dos modelos, que, no caso, foi o participativo.
No que diz respeito ao item B, que se referia aos três elementos do processo comunicacional, a maioria dos estudantes não conseguiu apresentar uma explicação clara e objetiva sobre o assunto.
Embora alguns tenham conseguido citar os elementos corretamente, eles se limitaram a isso e não foram capazes de desenvolver suas respostas de forma adequada.
Por outro lado, houve aqueles que tentaram explicar, de forma mais detalhada, como esse processo ocorre no modelo de administração participativa.
No entanto, demonstraram pouco conhecimento em relação às regras gramaticais e à estrutura da Língua Portuguesa.
Além disso, as respostas fornecidas pelos estudantes aos dois itens apresentaram informações que não estavam relacionadas ao contexto da questão.
Alguns citaram exemplos de gestão e organização empresarial, o que não foi solicitado na questão.
Essa falta de coerência entre o enunciado da questão e as respostas fornecidas evidencia uma lacuna na compreensão dos estudantes em relação ao conteúdo abordado.
Embora tenham demonstrado uma boa noção do modelo diretivo e participativo e tenham conseguido citar os fundamentos básicos, os estudantes enfrentaram dificuldades na resposta ao item B, que exigia a citação de três elementos do processo comunicacional.
As respostas que foram consideradas fracas foram aquelas às quais os estudantes responderam erroneamente ou parcialmente a todos os itens da questão.
Já as respostas medianas, no item A, abordavam o modelo participativo, de forma superficial, ou citavam apenas um fundamento.
No caso do item B, os estudantes não souberam citar claramente os três elementos do processo comunicacional.
Por outro lado, as respostas consideradas boas foram aquelas em que os estudantes conseguiram responder completamente ou quase completamente aos itens solicitados na questão.
Isso demonstra que alguns estudantes tiveram um bom entendimento do conteúdo abordado na questão e conseguiram aplicá -lo corretamente nas suas respostas.
Ficou evidente que os estudantes tinham o conhecimento dos fundamentos dos modelos, porém houve dificuldade em descrever os três elementos do processo comunicacional e associá -los ao modelo participativo relacionado ao comando do item b. 7.3.2.5 ANÁLISE DA QUESTÃO DISCURSIVA 5 DO COMPONENTE DE CONHECIMENTO ESPECÍFICO Na Tabela 7.20, constam as informações relativas à Questão 5 do conjunto do Componente de Conhecimento Específico.
Como revelam os dados da Tabela 7.20, o desempenho dos estudantes de todo o Brasil na Questão Discursiva 5 do Componente de Conhecimento Específico foi o 3.º maior entre as demais questões discursivas.
A nota média dos estudantes de todo o Brasil foi 30.
A maior Média foi registrada na região Sudeste ( 31,1 ), enquanto a menor Média foi registrada na região Nordeste ( 25,4 ).
Quanto à variabilidade das notas, o desvio-padrão dos estudantes do Brasil, como um todo, foi 32,6.
Enquanto o maior desvio foi encontrado na região Sudeste ( 33 ), o menor foi encontrado na região Centro-Oeste ( 30,7 ).
A mediana para o Brasil e para maioria das regiões foi 20, com exceção da região Nordeste.
Para o conjunto de estudantes de Tecnologia em Processos Gerenciais do Brasil, a nota máxima foi 100,0, obtida por pelo menos um estudante em cada região.
A nota mínima do Brasil foi zero, o mesmo valor encontrado em todas as regiões.
No Gráfico 7.13, é representada a distribuição das notas da Questão Discursiva 5 do Componente de Conhecimento Específico.
A distribuição mostrada no Gráfico 7.13 tem moda no intervalo ( 70 ; 80 ], com mais de 10 \% do total, excetuando -se a não pontuação, que foi mais de 40 \% dos estudantes que deixaram a questão em branco. 7.3.2.6 COMENTÁRIOS SOBRE AS RESPOSTAS À QUESTÃO DISCURSIVA 5 A questão abordou um \textbf{panorama} geral no que tange ao ponto crucial dentro na Administração de Empresas : o Planejamento Estratégico.
Na sua divisão orgânica, esse elemento fundamental traz alguns elementos que foram solicitados na questão, são eles : Missão, Visão, Mercado Potencial, Concorrentes e Consumidores.
Era dever do estudante, para atingir a nota máxima, relacionar tais pontos à elaboração de um plano de negócios – essa é outra ferramenta gerencial que elabora um escopo geral de um negócio em fase de ideação ou implementação estratégica dentro de um mercado.
Importante, também, notar que, ao pensar nos elementos do Planejamento Estratégico que a questão solicitou para lançamentos de respostas, é possível estabelecer noções de responsabilidade social de uma empresa e como sua vida orgânica está alinhada com tal proposição.
Missão e Visão se relacionam diretamente com o meio social que aquela instituição está envolvida e no qual seus processos são realizados.
Em fase relacional, observou -se, por inúmeras vezes, o deslocamento das proposições ideais de respostas ( padrão ) e das efetivas escritas pelos alunos.
O padrão resposta estabelece que a Missão é o objetivo de existência de uma empresa, isto é, para que ela foi criada ou aquilo que ela oferece como seus produtos e serviços.
Já a Visão é a objetificação do futuro da organização, ou seja, é preciso traçar uma meta de algo ou alguma coisa que a empresa queira alcançar no futuro, a partir de parâmetros bem definidos e já em fase de execução para alcançar o esperado.
Para esse ponto da questão, mesmo havendo respostas aceitas de maneira correta, inúmeros estudantes não conseguiram alcançar êxito por vários motivos, um dos mais importantes foi a troca de associação entre conceitos.
Isso foi \textbf{visualizado} ao encontrar respostas que, no lugar do conceito ou exemplificação da Missão, os discentes colocaram respostas sobre a Visão e vice-versa.
No que se refere ao ponto dois, era necessário que os candidatos possuíssem uma visão dos shareholder ( agentes envolvidos ) bem definida.
O mercado potencial deveria ser minuciosamente analisado, ou seja, o plano de negócios deveria estar com uma clara noção de possíveis segmentos ou nichos de clientes e mercado que a empresa poderia explorar.
Já os concorrentes e os consumidores são agentes fundamentais no dia a dia da criação de estratégias.
Também para esses casos, houve vários resultados positivos, inclusive, bem mais que o item anterior, contudo, ainda foi possível verificar dificuldade em apresentar pensamento analítico contextualizado aplicado aos negócios, isto é, o alinhamento entre estratégias e contexto social ou humano ao redor das empresas era difícil de se criar ou até ausente nas respostas obtidas.
Alguns pontos devem ser observados.
Embora o nível da questão seja considerado básico ou inicial, haja vista que a temática proposta é assunto de alicerce para outras disciplinas dentro da área, Há, também, uma apreciação que necessita ser feita : a caligrafia dos candidatos.
Em vários casos, era irreconhecível ou inidentificável a combinação de letras, frases e orações.
A semântica, os padrões ortográficos e sintáticos também são elementos considerados como negativos, ausentes ou com baixa presença em diversas respostas, isso evidencia que o conhecimento da Língua Portuguesa foi “ fraco ” entre os candidatos.
Outro fator preocupante é, novamente, a vasta quantidade de respostas “ em branco ” da questão, isso revela que, para esses casos, ocorreu a ausência total de aprendizagem dos elementos de estudo solicitados na prova.
Agora, voltando -se mais fundamentalmente às substâncias das respostas, alguns padrões foram encontrados.
O primeiro item solicitava que o discente tivesse afinidade com os conceitos de Missão e Visão.
Como evidenciado acima, Missão é o propósito existencial de uma organização e Visão representa seus objetivos futuros.
Um detalhe para esse caso precisa ser observado : a utilização e a conjugação de verbos.
Foi o caso de algumas respostas de alguns alunos que, ao responderem o que era a Missão da empresa, colocaram verbos na conjugação do futuro, isto é, indicando um pressuposto da Visão e não da Missão, deixando, assim, a questão sem resposta correta.
Além disso, foi fácil notar inúmeros casos em que os dois contextos eram trocados um pelo outro.
Sobre o segundo item de ordenação pertencente ao discriminatório de Mercado Potencial, Concorrentes e Consumidores, as respostas foram bem mais satisfatórias que o primeiro item ( Missão e Visão ).
Vários alunos conseguiram identificar os concorrentes e consumidores, bem como os conceitos ligados a esses termos.
No que se refere ao Mercado Potencial, houve várias respostas incorretas.
Aqui, destaca -se uma maior dificuldade entre os alunos para encontrar potenciais e oportunos mercados, seus segmentos e os nichos proveitosos em que uma organização pode trabalhar.
Por fim, foram consideradas como respostas regulares ou corretas aquelas em que os candidatos traçaram ideias claras, concisas e analíticas.
As relações com o contexto social, os agentes envolvidos e suas comunidades também foram pontos positivos.
Pensamento reflexivo, positivo e livre de amarras ideológicas sobre as entidades empresariais também foram considerados como corretos.
O padrão de resposta, por conceituação ou por exemplificação, também foi aceito para fins de avaliação e pontuação.
O corpo da questão e suas indagações solicitava aprendizagens básicas da área.
Contudo, mesmo assim, inúmeros candidatos apresentaram baixo ou nenhum conhecimento apreendido sobre tal.
Os ainda assim as respostas encontradas estavam na variância abaixo da média.
Foi possível observar a baixa capacidade de pensamento estratégico, lançamento de estratégias, noções espaciais e temporais, relacionamento entre agentes e atividades ( shareholder ) de um negócio, entre outros pontos.
Há, também, uma apreciação que necessita ser feita : a caligrafia dos candidatos.
Em vários casos, era irreconhecível ou inidentificável a combinação de letras, frases e orações.
A semântica, os padrões ortográficos e sintáticos também são elementos considerados como negativos, ausentes ou com baixa presença em diversas respostas, isso evidencia que o conhecimento da Língua Portuguesa foi “ fraco ” entre os candidatos.
Outro fator preocupante é, novamente, a vasta quantidade de respostas “ em branco ” da questão, isso revela que, para esses casos, ocorreu a ausência total de aprendizagem dos elementos de estudo solicitados na prova.
Agora, voltando -se mais fundamentalmente às substâncias das respostas, alguns padrões foram encontrados.
O primeiro item solicitava que o discente tivesse afinidade com os conceitos de Missão e Visão.
Como evidenciado acima, Missão é o propósito existencial de uma organização e Visão representa seus objetivos futuros.
Um detalhe para esse caso precisa ser observado : a utilização e a conjugação de verbos.
Foi o caso de algumas respostas de alguns alunos que, ao responderem o que era a Missão da empresa, colocaram verbos na conjugação do futuro, isto é, indicando um pressuposto da Visão e não da Missão, deixando, assim, a questão sem resposta correta.
Além disso, foi fácil notar inúmeros casos em que os dois contextos eram trocados um pelo outro.
Sobre o segundo item de ordenação pertencente ao discriminatório de Mercado Potencial, Concorrentes e Consumidores, as respostas foram bem mais satisfatórias que o primeiro item ( Missão e Visão ).
Vários alunos conseguiram identificar os concorrentes e consumidores, bem como os conceitos ligados a esses termos.
No que se refere ao Mercado Potencial, houve várias respostas incorretas.
Aqui, destaca -se uma maior dificuldade entre os alunos para encontrar potenciais e oportunos mercados, seus segmentos e os nichos proveitosos em que uma organização pode trabalhar.
Por fim, foram consideradas como respostas regulares ou corretas aquelas em que os candidatos traçaram ideias claras, concisas e analíticas.
As relações com o contexto social, os agentes envolvidos e suas comunidades também foram pontos positivos.
Pensamento reflexivo, positivo e livre de amarras ideológicas sobre as entidades empresariais também foram considerados como corretos.
O padrão de resposta, por conceituação ou por exemplificação, também foi aceito para fins de avaliação e pontuação.
O corpo da questão e suas indagações solicitava aprendizagens básicas da área.
Contudo, mesmo assim, inúmeros candidatos apresentaram baixo ou nenhum conhecimento apreendido sobre tal.
Os Os comportamentos de escrita também prejudicaram algumas notas pela análise irreconhecível de termos, palavras, frases, orações e outras partes.
Conhecimentos prévios e elementares, como os da própria Língua Portuguesa, foram ignorados ou esquecidos por diversos candidatos.
Observou -se, ainda, pouca capacidade de relacionamento entre o contexto social, humano, político e cultural do meio que a empresa está inserida e de suas funções e processos de trabalho ou de transformação em serviços e produtos.
Das respostas consideradas como acerto máximo, notou -se a rápida associação de ideias, a correta captação de conceitos, a exemplificação de objetos de maneira fiel e boa linguagem escrita dos termos solicitados e deferidos pela questão do Exame em questão. 7.3.2.7 CONSIDERAÇÕES FINAIS SOBRE AS QUESTÕES DISCURSIVAS DO COMPONENTE DE CONHECIMENTO ESPECÍFICO As respostas encontradas parecem revelar uma área de formação pensada na sociedade e suas aplicações concretas, suas demandas existentes e atualizadas pela nova globalização e pelo contexto tecnológico atual.
Embora tenham ocorrido inúmeros erros nas respostas, ainda assim o padrão de gerenciamento de pessoas e processos foi encontrado ora explícito ora implícito.

% matched lemmas: panorama, visualizar
\end{textsample}
